\documentclass[aps,prd,reprint,onecolumn,nofootinbib]{revtex4-2}
\usepackage[T1]{fontenc}
\usepackage[utf8]{inputenc}
\usepackage{amsmath,amssymb,graphicx,booktabs}
\usepackage{xurl}\usepackage[colorlinks=true]{hyperref}
\setlength{\emergencystretch}{4em}
\InputIfFileExists{generated/numbers.tex}{}{\PackageError{PR4}{PENDING_SOURCE_QUALIFICATION}{Final qualified MC100 source required.}\stop}
\begin{document}
\title{Low-multipole polarization morphology in processed Commander maps}
\author{Conditional methods research draft}\noaffiliation
\maketitle
\hypertarget{observable-and-processing}{%
\section{Observable and processing}\label{observable-and-processing}}

We compare a fixed polarization morphology statistic in the Planck PR4
Commander full-mission map with the corresponding released
NPIPE/Commander simulations. This is an experiment on processed maps.
The related finite-bandpass calculation uses the 24 real signed
coordinates of E and B at multipoles 2 and 3. Here multipole 4 is
additionally required by the statistic. The experiment supplies neither
a new spectral-rank result nor a measured bound on foreground moments.

The observed polarization comes from the NSIDE = 2048 Q and U extension
of \nolinkurl{COM_CMB_IQU-commander_4096_R4.00_full.fits}; the
higher-resolution temperature extension is unused. Simulations are the
fixed 100 full-mission products
\nolinkurl{npipe6v20_commander_cmb_2048_mc_0300_005a.fits} through
\nolinkurl{..._0399_005a.fits}. Each simulation identifier contributes
one parent-sky realization. Temporal splits and alternative products of
the same sky do not add independent realizations.

We retain Galactic coordinates and the released COSMO polarization
convention. Observed Q and U are in K\_CMB; simulated values are
converted from microkelvin to kelvin before analysis. The binary
Lollipop mask retains 1273 of 3072 NSIDE = 16 pixels. Each NESTED parent
value is copied to its 16384 NSIDE = 2048 children, and the same mask is
applied to native-resolution Q and U. The masked maps are reordered to
RING for spin-two analysis with lmax = mmax = 4:

\[
 \widetilde a^X_{\ell m}=\mathcal A^X_{\ell m}[M Q,M U],
 \qquad X\in\{E,B\},\quad \ell=2,3,4.
\]

We use the delivered maps without additional smoothing, beam inversion
or noise subtraction. No scalar effective beam is specified or inferred;
the delivered processing response and the additional mask response are
retained. The complete signed E and B coefficients, satisfying
\(\widetilde a^X_{\ell,-m}=(-1)^m\overline{\widetilde a^X_{\ell m}}\),
are retained with their source, unit and processing metadata.

\hypertarget{a-fixed-parity-odd-contraction}{%
\section{A fixed parity-odd
contraction}\label{a-fixed-parity-odd-contraction}}

For each nonzero E multipole define

\[
 v_{\ell m}=\frac{\widetilde a^E_{\ell m}}
 {\left(\sum_{r=-\ell}^{\ell}|\widetilde a^E_{\ell r}|^2\right)^{1/2}},
 \quad
 J_{234}=\operatorname{Im}\sum_{m+n+p=0}
 \begin{pmatrix}2&3&4\\m&n&p\end{pmatrix}
 v_{2m}v_{3n}v_{4p},\qquad T=|J_{234}|.
\]

The sum uses every allowed signed index, with no additional
symmetrization factor. Reality of the field and the Wigner-symbol
symmetry make the unprojected contraction purely imaginary: its
conjugate is \((-1)^{2+3+4}\) times itself. Common proper rotations
preserve the contraction; coefficient parity reverses J and preserves T.
These algebraic properties concern the extracted coefficients. Rotating
a sky through a fixed Galactic mask is a different operation.

T is dimensionless and uses one field, one multipole triple, one map
product and one mask. There is no maximization over these choices.
Powers
\(\widetilde C_\ell^X=(2\ell+1)^{-1}\sum_m|\widetilde a^X_{\ell m}|^2\),
in K\_CMB squared, and quadrupole--octupole axis alignment are
descriptive quantities. An undefined descriptive axis does not
invalidate an otherwise valid J. This contraction is distinct from
\(\chi=\det(v,Qv,Q^2v)\); the sharp bound for that determinant and ideal
spherical reference laws are not used to calibrate T.

\hypertarget{released-pipeline-reference-and-finite-simulation-precision}{%
\section{Released-pipeline reference and finite-simulation
precision}\label{released-pipeline-reference-and-finite-simulation-precision}}

The released simulation construction uses independent CMB realizations
with end-to-end processing, fixed foreground and instrumental settings,
and Commander mixing matrices estimated from the observed data, with
amplitudes fitted for each simulation. We use this as a plug-in
processed-map working law. Fixing estimated mixing matrices does not
establish the exact conditional distribution given their estimates or
unconditional physical coverage. The data products and processing
genealogy are documented by the NPIPE release and the Planck Legacy
Archive\textasciitilde{}\cite{NPIPE2020,ESA_SIMULATION,PLA_PRODUCTS}.

For the complete prescribed ensemble, define

\[
 B=100,\qquad k=\sum_{b=1}^{100}{\bf1}\{T_b\ge T_{\rm obs}\},
 \qquad \widehat\theta=k/B,\qquad p_{\rm rank}=(1+k)/(B+1).
\]

Ties are included. The smallest rank value is 1/101. Interpreting this
rank as a Monte Carlo significance reference additionally assumes
exchangeability of the observed statistic and simulated statistics under
the stipulated working model. No selection or early stopping uses
intermediate tail counts.

For the working-law tail probability
\(\theta=\Pr_{\rm work}(T_*\ge T_{\rm obs})\), we report the two-sided
95\% Clopper--Pearson interval from k successes in B independent
simulation draws. Its endpoints are 0 for k = 0 or
\(F^{-1}_{\mathrm{Beta}(k,B-k+1)}(0.025)\) otherwise, and 1 for k = B or
\(F^{-1}_{\mathrm{Beta}(k+1,B-k)}(0.975)\) otherwise. The interval uses
B binomial trials, not B + 1 pseudo-observations. Its half-width must
meet the specified 0.105 precision criterion. It measures
finite-simulation uncertainty in this tail probability, without
accounting for unrepresented foreground, calibration or model
uncertainty.

\hypertarget{results-and-interpretation}{%
\section{Results and interpretation}\label{results-and-interpretation}}

For all 100 prescribed realizations, the observed value is
\PRFourTobs{}, with \PRFourK{} inclusive exceedances. The empirical tail
fraction is \PRFourEmpirical{}, and the plus-one rank is \PRFourRank{}.
The 95\% binomial interval is \([\PRFourCPlow,\PRFourCPhigh]\), with
half-width \PRFourCPhalf{}. Figure 1 displays the complete simulation
distribution and observed statistic. The accompanying numerical tables
retain every simulation and its signed E/B coefficients.

Three unit-norm harmonic inputs, E(2,0), a real-field
E(3,\textbar m\textbar=1) mode and a real-field
B(4,\textbar m\textbar=2) mode, illustrate the native-resolution
mask-and-analysis response in Figure 2. Their outputs retain every E/B
component through multipole 4. They are three selected response columns;
they do not reconstruct the full analysis operator or an upstream NPIPE
transfer matrix. The qualified numerical supplement records linearity,
masked-pixel, unit, reality, parity, rotation and retained-coefficient
stability checks.

The comparison quantifies compatibility with this particular
processed-map working law. It does not measure a physical foreground
budget, infer radiation stress or matter shear from E, identify a
Bianchi family, or establish unconditional false-positive control.

\input{generated/methods_table.tex}
\input{generated/results_table.tex}

\begin{figure}[htbp]
\centering\includegraphics[width=.95\linewidth,height=.72\textheight,keepaspectratio]{figures/morphology_survival.pdf}
\caption{Empirical survival fraction of T for all 100 fixed processed Commander realizations, with the observed value marked. At the observed threshold the inclusive fraction is k/100; the separately reported rank is (1+k)/101. The interval in the result table describes simulation-tail uncertainty under the working law.}
\end{figure}

\begin{figure}[htbp]
\centering\includegraphics[width=.95\linewidth,height=.72\textheight,keepaspectratio]{figures/mask_response_columns.pdf}
\caption{Three selected columns of the common mask-and-spin-analysis response, shown in the real signed E/B representation at multipoles 2, 3 and 4. Rows within each multipole are m = 0 followed by (sqrt(2) Re, -sqrt(2) Im) for positive m. Each input has unit harmonic norm. Positive-m inputs for the last two columns are 1/sqrt(2) and i/sqrt(2), respectively, with negative m fixed by reality. The display includes inter-multipole and E/B coupling where present; it is not a complete transfer matrix.}
\end{figure}

\bibliographystyle{apsrev4-2}\bibliography{references}
\end{document}
